\begin{question}[ID=1.1, subtitle={Calcolo di derivate — serie 1}, tags={derivate, calcolo}, difficulty=media]
Calcola la derivata delle seguenti funzioni, semplificando dove possibile:
\begin{center}\begin{minipage}{0.92\linewidth}
\begin{multicols}{2}
\begin{enumerate}[label=\arabic*),leftmargin=3.2cm]
\item $f(x) = x^3 - 5x^2 + 2x - 7$
\item $f(x) = \sqrt{x^2 + 4x + 4}$
\item $f(x) = \frac{2x}{x^2 + 1}$
\item $f(x) = x^3 \cdot e^x$
\item $f(x) = \ln(x^2 + 1)$
\item $f(x) = \arcsin(x-1)$
\item $f(x) = x \cdot \tan x$
\item $f(x) = 2^{-x} \cdot \cos(x)$
\item $f(x) = \frac{|x - 1|}{x}$
\item $f(x) = x^{3x}$
\end{enumerate}
\end{multicols}
\end{minipage}\end{center}
\begin{solution}
1) $f'(x)=3x^2-10x+2$. \\
2) $\sqrt{x^2+4x+4}=|x+2|$, quindi $f'(x)=\dfrac{x+2}{|x+2|}$ ($x\neq-2$). \\
3) $f'(x)=\dfrac{2(1-x^2)}{(x^2+1)^2}$. \\
4) $f'(x)=x^2e^x(3+x)$. \\
5) $f'(x)=\dfrac{2x}{x^2+1}$. \\
6) $f'(x)=\dfrac{1}{\sqrt{1-(x-1)^2}}=\dfrac{1}{\sqrt{2x-x^2}}$ ($0<x<2$). \\
7) $f'(x)=\tan x+\dfrac{x}{\cos^2 x}$. \\
8) $f'(x)=-2^{-x}(\ln 2\,\cos x+\sin x)$. \\
9) $\dfrac{|x-1|}{x}$: $f'(x)=\dfrac{1}{x^2}$ per $x>1$, $f'(x)=-\dfrac{1}{x^2}$ per $x<1$; non derivabile in $x=1$ (punto angoloso). \\
10) $x^{3x}=e^{3x\ln x}$, quindi $f'(x)=3x^{3x}(\ln x+1)$.
\end{solution}
\end{question}

%-----------------------------------------------------%

\begin{question}[ID=1.2, subtitle={Calcolo di derivate — serie 2}, tags={derivate, calcolo}, difficulty=media]
Calcola la derivata delle seguenti funzioni, semplificando dove possibile:
\begin{center}\begin{minipage}{0.92\linewidth}
\begin{multicols}{2}
\begin{enumerate}[label=\arabic*),leftmargin=3.2cm]
\item $f(x) = x^4 - 3x^3 + 2x^2 - 1$
\item $f(x) = \sqrt{1 + x^4}$
\item $f(x) = \frac{1}{\sqrt{x^2 + 1}}$
\item $f(x) = x^2 \cdot \ln(x)$
\item $f(x) = x^{2x}$
\item $f(x) = \cos(x^3)$
\item $f(x) = \frac{x+2}{x^2 - 1}$
\item $f(x) = e^x \cdot \sin(x)$
\item $f(x) = |x^2 - 4|$
\item $f(x) = \ln(\sqrt{x^2 + 1})$
\end{enumerate}
\end{multicols}
\end{minipage}\end{center}
\begin{solution}
1) $f'(x)=4x^3-9x^2+4x$. \\
2) $f'(x)=\dfrac{2x^3}{\sqrt{1+x^4}}$. \\
3) $f'(x)=-\dfrac{x}{(x^2+1)^{3/2}}$. \\
4) $f'(x)=x(2\ln x+1)$. \\
5) $x^{2x}=e^{2x\ln x}$, quindi $f'(x)=2x^{2x}(\ln x+1)$. \\
6) $f'(x)=-3x^2\sin(x^3)$. \\
7) $f'(x)=\dfrac{-(x^2+4x+1)}{(x^2-1)^2}$. \\
8) $f'(x)=e^x(\sin x+\cos x)$. \\
9) $|x^2-4|$: $f'(x)=2x$ per $|x|>2$, $f'(x)=-2x$ per $|x|<2$; non derivabile in $x=\pm2$. \\
10) $\ln\sqrt{x^2+1}=\tfrac12\ln(x^2+1)$, quindi $f'(x)=\dfrac{x}{x^2+1}$.
\end{solution}
\end{question}

%-----------------------------------------------------%

\begin{question}[ID=2.1, subtitle={Retta tangente in tre punti assegnati}, tags={derivate, retta tangente}, difficulty=media]
Determina l'equazione della retta tangente alla curva data nei seguenti punti:
\begin{center}\begin{minipage}{0.85\linewidth}
\begin{enumerate}[label=\arabic*)]
\item $f(x) = x^2 + 2x - 3, \quad x_0 = -1$
\item $f(x) = e^{x^2}, \quad x_0 = 0$
\item $f(x) = \ln(x+1), \quad x_0 = 1$
\end{enumerate}
\end{minipage}\end{center}
\begin{solution}1) $y=-4$; 2) $y=1$; 3) $y=\tfrac{1}{2}(x-1)+\ln 2$\end{solution}
\end{question}

%-----------------------------------------------------%

\begin{question}[ID=2.2, subtitle={Retta tangente con verifica intersezione asse $y$}, tags={derivate, retta tangente}, difficulty=media]
Trova l'equazione della retta tangente alla curva nei seguenti punti e verifica se la tangente interseca l'asse \( y \):
\begin{center}\begin{minipage}{0.85\linewidth}
\begin{enumerate}[label=\arabic*)]
\item $f(x) = x^3 - 2x + 2, \quad x_0 = 1$
\item $f(x) = \cos(x), \quad x_0 = \frac{\pi}{3}$
\item $f(x) = \sqrt{x+2}, \quad x_0 = 2$
\end{enumerate}
\end{minipage}\end{center}
\begin{solution}1) $y=x$; 2) $y=-\tfrac{\sqrt{3}}{2}(x-\tfrac{\pi}{3})+\tfrac{1}{2}$; 3) $y=\tfrac{1}{4}(x-2)+2$\end{solution}
\end{question}

%-----------------------------------------------------%

\begin{question}[ID=3.1, subtitle={Studio di $f(x) = x^3/(x^2-4)$ con derivate}, tags={derivate, studio di funzione}, difficulty=difficile]
Studia la seguente funzione, facendo particolare attenzione allo studio del segno della derivata prima e seconda:
\begin{equation*}
    f(x) = \frac{x^3}{x^2-4}
\end{equation*}
\begin{solution}
\textbf{Dominio e simmetrie:} $x\neq\pm2$; funzione dispari ($f(-x)=-f(x)$). \\
\textbf{Asintoti:} verticali $x=\pm2$. Dalla divisione $\dfrac{x^3}{x^2-4}=x+\dfrac{4x}{x^2-4}$ segue l'asintoto obliquo $y=x$. \\
\textbf{Derivata prima:} $f'(x)=\dfrac{x^2(x^2-12)}{(x^2-4)^2}$. Poiché $x^2\ge0$ e il denominatore è positivo, $f'>0$ per $|x|>2\sqrt3$ e $f'<0$ per $|x|<2\sqrt3$ (escluso $x=\pm2$). Quindi massimo in $x=-2\sqrt3$ ($f=-3\sqrt3$), minimo in $x=2\sqrt3$ ($f=3\sqrt3$); in $x=0$ tangente orizzontale senza estremo (flesso). \\
\textbf{Derivata seconda:} $f''(x)=\dfrac{8x(x^2+12)}{(x^2-4)^3}$, che si annulla in $x=0$: flesso in $(0,0)$. Concavità verso l'alto per $-2<x<0$ e $x>2$, verso il basso per $x<-2$ e $0<x<2$.
\end{solution}
\end{question}

%-----------------------------------------------------%

\begin{question}[ID=3.2, subtitle={Studio di $f(x) = x^3/(x-4)$ con derivate}, tags={derivate, studio di funzione}, difficulty=difficile]
Studia la seguente funzione, facendo particolare attenzione allo studio del segno della derivata prima e seconda:
\begin{equation*}
    f(x) = \frac{x^3}{x- 4}
\end{equation*}
\begin{solution}
\textbf{Dominio:} $x\neq4$. \\
\textbf{Asintoti:} verticale $x=4$. Dalla divisione $\dfrac{x^3}{x-4}=x^2+4x+16+\dfrac{64}{x-4}$ si vede che non c'è asintoto obliquo, ma un asintoto \emph{parabolico} $y=x^2+4x+16$. \\
\textbf{Derivata prima:} $f'(x)=\dfrac{2x^2(x-6)}{(x-4)^2}$. Essendo $x^2\ge0$ e denominatore positivo, $f'>0$ per $x>6$ e $f'<0$ per $x<6$ (escluso $x=4$); minimo in $x=6$ con $f(6)=108$. In $x=0$ tangente orizzontale senza estremo. \\
\textbf{Derivata seconda:} $f''(x)=2+\dfrac{128}{(x-4)^3}$, che si annulla in $x=0$: flesso in $(0,0)$. Concavità verso l'alto per $x<0$ e $x>4$, verso il basso per $0<x<4$.
\end{solution}
\end{question}

%-----------------------------------------------------%

\begin{question}[ID=4.1, subtitle={Segno della derivata di $ke^{kx}-e^{-2kx}$ in $x=0$}, tags={derivate, parametrico}, difficulty=media]
Determina per quali valori di $k$ la derivata della funzione $f(x) = ke^{kx} - e^{-2kx}$ nel punto $x=0$ risulti positiva.
\begin{solution}$k^2+2k>0 \Rightarrow k<-2 \vee k>0$\end{solution}
\end{question}

%---------------------------------------------------%

\begin{question}[ID=4.2, subtitle={Derivata di $k\ln(x+2)+e^{-kx}$ in $x=1$ maggiore o uguale a 1}, tags={derivate, parametrico}, difficulty=media]
Determina per quali valori di \( k \) la derivata di $f(x)$ nel punto \( x = 1 \) è maggiore o uguale a 1.
\[
f(x) = k\ln(x+2) + e^{-kx}
\]
\begin{solution}
$f'(x)=\dfrac{k}{x+2}-ke^{-kx}$. In $x=1$: $f'(1)=\dfrac{k}{3}-ke^{-k}=k\!\left(\dfrac13-e^{-k}\right)$. La condizione $f'(1)\ge1$ è trascendente; risolta numericamente dà $k\lesssim-0{,}70$ oppure $k\gtrsim3{,}45$.
\end{solution}
\end{question}

%---------------------------------------------------%

\begin{question}[ID=4.3, subtitle={Derivata di $\sin(kx)+\cos(2kx)$ positiva in $x=\pi/6$}, tags={derivate, parametrico}, difficulty=media]
Per la funzione
\[
f(x) = \sin(kx) + \cos(2kx),
\]
calcola la derivata e determina i valori di \( k \) tali che la derivata calcolata nel punto \( x = \dfrac{\pi}{6} \) sia positiva.
\begin{solution}
$f'(x)=k\cos(kx)-2k\sin(2kx)=k\big[\cos(kx)-2\sin(2kx)\big]$. In $x=\dfrac{\pi}{6}$: $f'\!\left(\dfrac{\pi}{6}\right)=k\Big[\cos\dfrac{k\pi}{6}-2\sin\dfrac{k\pi}{3}\Big]>0$. La parentesi dipende da $k$ in modo non elementare (funzione oscillante), quindi la disequazione si risolve graficamente: il segno di $f'$ è quello di $k$ moltiplicato per quello della parentesi.
\end{solution}
\end{question}

%---------------------------------------------------%

\begin{question}[ID=4.4, subtitle={Derivata di $kx/(x+1)-e^{-kx}$ negativa in $x=2$}, tags={derivate, parametrico}, difficulty=media]
Considera la funzione
\[
f(x) = \frac{kx}{x+1} - e^{-kx}.
\]
Calcola la derivata e trova i valori di \( k \) per i quali la derivata nel punto \( x = 2 \) è negativa.
\begin{solution}
$f'(x)=\dfrac{k}{(x+1)^2}+ke^{-kx}=k\!\left[\dfrac{1}{(x+1)^2}+e^{-kx}\right]$. In $x=2$: $f'(2)=k\!\left(\dfrac19+e^{-2k}\right)$. La parentesi è sempre positiva, quindi $f'(2)<0\iff k<0$.
\end{solution}
\end{question}

%---------------------------------------------------%

\begin{question}[ID=4.5, subtitle={Derivata di $kx^2+e^{kx}-\ln(x+1)$ in $x=0$ maggiore o uguale a 2}, tags={derivate, parametrico}, difficulty=media]
Per la funzione
\[
f(x) = kx^2 + e^{kx} - \ln(x+1),
\]
calcola la derivata e determina i valori di \( k \) tali che la derivata calcolata nel punto \( x = 0 \) sia maggiore o uguale a 2.
\begin{solution}$f'(0) = k + 1 \ge 2 \Rightarrow k \ge 1$\end{solution}
\end{question}
